\documentclass{beamer}

\mode<presentation>
{
  \usetheme{Singapore}      % or try Darmstadt, Madrid, Warsaw, ...
  \usecolortheme{crane} % or try albatross, beaver, crane, ...
  \usefonttheme{serif}  % or try serif, structurebold, ...
  \setbeamertemplate{navigation symbols}{}
  \setbeamertemplate{caption}[numbered]
} 

\usepackage[autostyle=true]{csquotes}
\usepackage{xeCJK}
\usepackage[T1]{fontenc}
\usepackage{xcolor}
\usepackage{setspace}
\setmainfont[AutoFakeSlant=0.2]{Charis SIL}
\setCJKmainfont[AutoFakeSlant=0.2]{Noto Serif CJK HK}
\setCJKmonofont[AutoFakeSlant=0.2]{Noto Serif CJK HK}
\usepackage{polyglossia}
\setmainlanguage{english}
\usepackage{csquotes}
\usepackage{array}
\usepackage{makecell}
\usepackage{changepage}

\title[Pearl River Sinitic]{Sinitic Languages of the Pearl River Region \\ --- Unit 3.1}
\subtitle{13ᵗʰ EACL Summer School, Universität Stuttgart}
\author{Hilário de Sousa 蘇沙}
\institute{EHESS–CRLAO \\ {\tiny membre associé}}
\date{Wednesday 29 July 2026 \\ 丙午六月十六 \\ {\footnotesize \texttt{hilario@bambooradical.com}}}

\begin{document}

\begin{frame}
  \titlepage
\end{frame}

\section{MC finals}

\begin{frame}[fragile]{MC finals}
\textbf{攝[摄] sip3/shè final groups}:
\begin{table}
    \centering
    \renewcommand{\arraystretch}{1.2}
    \begin{tabular}{l l l l|l l l l}
        果 & kwa\textsuperscript{B} & gwo2 & guǒ & 咸 & ɦɛm\textsuperscript{A} & \textcolor{blue}{haam4} & xián\\
        假 & kæ\textsuperscript{B} & gaa2 & jiǎ & 深 & ɕim\textsuperscript{A} & \textcolor{red}{sam1} & shēn\\ \cline{5-8}
        遇 & ŋju\textsuperscript{C} & jyu6 & yù & 山 & ʂɛn\textsuperscript{A} & \textcolor{blue}{saan1} & shān\\ \cline{1-4}
        蟹 & ɦɛ\textsuperscript{B} & haai5 & xiè & 臻 & ʈʂin\textsuperscript{A} & \textcolor{red}{zeon1} & zhēn\\ \cline{5-8}
        止 & tɕi\textsuperscript{B} & zi2 & zhǐ & 宕 & daŋ\textsuperscript{C} & \textcolor{blue}{dong6} & dàng\\ \cline{1-4}
        效 & ɦæw\textsuperscript{C} & \textcolor{blue}{haau6} & xiào & 江 & kæʷŋ\textsuperscript{A} & \textcolor{blue}{gong1} & jiāng\\
        流 & ljuw\textsuperscript{A} & \textcolor{red}{lau4} & liú & 曾 & tsəŋ\textsuperscript{A} & \textcolor{red}{zang1} & zēng\\
         & & & & 梗 & kɐŋ\textsuperscript{B}& \textcolor{red}{gang2} & gěng\\
         & & & & 通 & tʰuʷŋ\textsuperscript{A} & \textcolor{red}{tung1} & tōng\\
    \end{tabular}
    \caption{Middle Chinese 攝 final groups}
\end{table}

\end{frame}

\begin{frame}[fragile]{MC finals}

Cantonese: more than 90\% of the time, characters in the final groups of
        \begin{itemize}
            \item \textcolor{blue}{效 haau6 咸 haam4 山 saan1 宕 dong6 江 gong1}\\ have a long vowel
            \item \textcolor{red}{流 lau4 深 sam1 臻 zeon1 曾 zang1 梗 gang2 通 tung1}\\ have a short vowel
        \end{itemize}

\end{frame}

\begin{frame}[fragile]{MC finals}

\begin{table}
    \centering
      {\scriptsize
        \setlength{\tabcolsep}{1pt}
        \renewcommand{\arraystretch}{1.2}    \begin{tabular}{c|c c c c|c c c c}
         & 開一 & 開二 & 開三 & 開四 & 合一 & 合二 & 合三 & 合四 \\ \hline
        果 & 歌 a & & 戈 ja & & 戈 wa & & 戈 wja & \\ \hline
        假 & & 麻 æ & 麻 jæ & & & 麻 wæ & & \\ \hline
        遇 & & & & & 模 u & & 魚 jə & \\
        遇 & & & & & & & 虞 ju & \\ \hline
        蟹 & 咍 əj & 皆 ɛj & 祭 j(i)ej & 齊 ej & 灰 wəj & 皆 wɛj & 祭 wj(i)ej & 齊 wej \\
        蟹 & 泰 aj & 佳 ɛ & 廢 jəj & & 泰 waj & 佳 wɛ & 廢 wjəj & \\
        蟹 & & 夬 æj & & & & 夬 wæj & & \\ \hline
        止 & & & 支 j(i)e & & & & 支 wj(i)e & \\
        止 & & & 脂 (j)ij & & & & 脂 w(j)ij & \\
        止 & & & 之 i & & & & & \\
        止 & & & 微 jɨj & & & & 微 wjɨj & \\ \hline
        效 & 豪 aw & 肴 æw & 宵 j(i)ew & 蕭 ew & & & & \\ \hline
        流 & 侯 uw & & 尤 juw & & & & & \\
        流 & & & 幽 jiw & & & & & \\  
    \end{tabular}}
\end{table}

\end{frame}

\begin{frame}[fragile]{MC finals}

\begin{table}
    \centering
      {\scriptsize
        \setlength{\tabcolsep}{1pt}
        \renewcommand{\arraystretch}{1.2}    \begin{tabular}{c|c c c c|c c c c}
         & 開一 & 開二 & 開三 & 開四 & 合一 & 合二 & 合三 & 合四 \\ \hline
        咸 & 覃合 əm & 咸洽 ɛm & 鹽葉 j(i)em & 添帖 em & & & 凡乏 wjəm & \\
        咸 & 談盍 am & 銜狎 æm & 嚴業 jæm & & & & & \\ \hline 
        深 & & & 侵緝 (j)im & & & & & \\ \hline
        山 & 寒曷 an & 山黠 ɛn & 仙薛 j(i)en & 先屑 en & 桓末 wan & 山黠 wɛn & 仙薛 wj(i)en & 先屑 wen \\
        山 & & 刪轄 æn & 元月 jən & & & 刪轄 wæn & 元月 wjən & \\ \hline
        臻 & 痕 ən & & 真質/臻櫛 (j)in & & 魂沒 wən & & 諄術 w(j)in & \\
        臻 & & & 殷迄 jɨn & & & & 文物 jun & \\ \hline
        宕 & 唐鐸 aŋ & & 陽藥 jaŋ & & 唐鐸 waŋ & & 陽藥 wjaŋ & \\ \hline
        江 & & 江覺 æʷŋ & & & & & & \\ \hline
        曾 & 登德 əŋ & & 蒸職 iŋ & & 登德 wəŋ & & 職 wik & \\ \hline
        梗 & & 庚陌 æŋ & 庚陌 jæŋ & 青錫 eŋ & & 庚陌 wæŋ & 庚 wjæŋ & 青 weŋ \\
        梗 & & 耕麥 ɛŋ & 清昔 j(i)eŋ & & & 耕麥 wɛŋ & 清昔 wj(i)eŋ & \\ \hline
        通 & & & & & 東屋 uʷŋ & & 東屋 juʷŋ & \\
        通 & & & & & 冬沃 əʷŋ & & 鐘燭 jəʷŋ & \\
    \end{tabular}}
\end{table}

\end{frame}

\begin{frame}[fragile]{MC finals}

MC codas: -j -w -m -p -n -t -ŋ -k -ʷŋ -ʷk\\
>95\% retained in Cantonese\\
{\footnotesize (-ʷŋ -ʷk as -ŋ -k)}\\
\medskip
vowels: all sorts of changes in Cantonese\\
{\footnotesize (very not conservative)}\\

\end{frame}

\section{Early Cantonese}

\begin{frame}[fragile]{Early Cantonese}
Earliest text with elements identifiable as Cantonese:\\
《花箋記》faa1 zin1 gei3\\ (late Ming: 17\textsuperscript{th} century, or perhaps 16\textsuperscript{th} century)
\begin{figure}
\centering
\includegraphics[width=0.55\textwidth]{faazingei2021}
{\tiny \url{https://upload.wikimedia.org/wikipedia/commons/a/a6/Faazin_Gei_20-21.jpg}}
\end{figure}

\end{frame}

\begin{frame}[fragile]{Early Cantonese}
Earliest Cantonese rime book:\\ 
《分韻撮要》fan1 wan5 cyut3 jiu3 (early Qing dynasty)\\
\smallskip
{\footnotesize Most-famous dated editions: 《江湖尺牘分韻撮要合集》(1782)\\《英華分韻撮要》 Tonic Dictionary of the Chinese Language in the Canton Dialect (1856)}\\
\medskip
Only recorded literary words / pronunciations\\
\medskip
\textit{\small (Many older traits of Standard Cantonese can still be found in many other Sinitic varieties in the wider region)}

\end{frame}

\begin{frame}[fragile]{Early Cantonese}
Sham 岑 (2020)'s reconstruction (from pre-1782 editions)\\
\begin{columns}
    \begin{column}{0.5\textwidth}
        \begin{figure}
        \includegraphics[width=1\textwidth]{fanwancyutjiuonsets}
        \end{figure}
    \end{column}
    \begin{column}{0.5\textwidth}
        \begin{itemize}
            \item /ts tsʰ s/ vs. /tɕ tɕʰ ɕ/\\
            鄒 zau1 秋 cau1 修 sau1\\
            周 žau1 抽 čau1 收 šau1\\
            (now no distinction)
            \item /ɲ/ vs. /j/\\
            魚 ñyu4 vs. 于[于] jyu4\\
            (now merged into /j/)\\
            (/ɲ/ < MC [ɲ] or [ŋ])
        \end{itemize}
    \end{column}
\end{columns}

\vfill \setstretch{0.5} {\tiny 
  Sham, Ricky Yiu Ho 岑堯昊. 2020.《分韻撮要》版本與音系新論 Innovations on Fenyun Cuoyao: Textual history and phonological reconstruction. \textit{Current Research in Chinese Linguistics} 中國語文通訊 99(1). \url{https://www.cuhk.edu.hk/ics/clrc/crcl_99_1/sham.pdf}\par} 
\end{frame}

\begin{frame}[fragile]{Early Cantonese}

\begin{columns}
    \begin{column}{0.8\textwidth}
        \begin{figure}
        \includegraphics[width=1\textwidth]{fanwancyutjiurimes}
        \end{figure}
    \end{column}
    \begin{column}{0.35\textwidth}
        {\footnotesize
        \begin{itemize}
            \item {[e o ɵ]} {\tiny ([ɪ ʊ ʏ]?)} >\\
            {[i u y]\textsuperscript{文} [ɛ ɔ œ]\textsuperscript{白}}\\
            \smallskip
            e.g. 誅伐 zyu1\textsuperscript{文} fat6\\ 
            zoe1\textsuperscript{白} zyu6 saai3\\
            \smallskip
            {(my POV: 苦}\\
            Canto. [fu˧˥]\\
            Nan. Ping. [hɔ˧]\\
            Viet. khổ\textsuperscript{文} / khó\textsuperscript{白})\\
            \item {[ɿ]([z̩]): allophone of [i] / ts tsʰ s\_}\\
            子 [tsz˧˥] 止 [tɕi˧˥]\\
            (cf. Brit. HK spelling of 子 Tsz) 
        \end{itemize}}
    \end{column}
\end{columns}
{\footnotesize
\begin{itemize}
     \item {[i y] > [ej ɵɥ] (except tɕ tɕʰ ɕ ȵ j\_):} \\ \underline{你 [nej˩˧]}, (子 [tsz˧˥],) 止 [tɕi˧˥], 耳 [ȵi˩˧], \underline{企(徛) [kʰej˩˧]}\\
     \underline{女 [nɵɥ˩˧]}, \underline{咀 [tsɵɥ˧˥]}, 主 [tɕy˧˥], 語 [ȵy˩˧], \underline{拒 [kʰɵɥ˩˧]}
\end{itemize}}
           
\end{frame}

\begin{frame}[fragile]{Early Cantonese}

{\footnotesize
\begin{itemize}
    \item {[u] > [ow] (except f kʷ kʷʰ w\_):}\\
    \underline{補 [pow˧˥]}, 甫 [fu˧˥], \underline{嫂 [sow˧˥]}, \underline{稿 [kow˧˥]}, 古 [kʷu˧˥]\\
    (cf. Brit. HK spelling of 古 Kwu; modern Canto. 稿 gou2, 古 gu2)\\
    \item In《分韻撮要》, there were no [ej ow], but [ɵɥ] already existed.\\
    {[y] > [ɵɥ] (except tɕ tɕʰ ɕ ȵ j\_)} \textsuperscript{($\approx$MC 魚 jə 虞 ju 遇合三)}\\
    \quad 慮 [ly˨] > [lɵɥ˨], 除 [tsʰy˩] > [tsʰɵɥ˩]\\
    {[ɵɥ] > [ɵɥ]} \textsuperscript{($\approx$MC 祭 wj(i)ej 蟹合三, 支 wj(i)e 脂 w(j)ij 止合三)}\\
    \quad 戾 [lɵɥ˨], 隨 [tsʰɵɥ˩]\\
    {[uj] > [ɵɥ] (only t tʰ ts tsʰ s n l\_)} \textsuperscript{($\approx$MC 灰 wəj 蟹合一)}\\
    \quad 雷 [luj˩] > [lɵɥ˩], 崔 [tsʰuj˥˧] > [tsʰɵɥ˥]\\
    \item {[ɔm ɔp] > [ɐm ɐp]} \textsuperscript{($\approx$MC 覃合 əm əp 談盍 am ap 咸開一)}\\ 
    窞 [tʰɔm˩˧] > [tʰɐm˩˧], 甘 [kɔm˥˧] > [kɐm˥],\\ 坎 [hɔm˧˥] > [hɐm˧˥], 暗 [ɔm˧] > [ɐm˧], 合 [hɔp˨] > [hɐp˨]\\
    (cf. Brit. HK spelling of 紅磡 Hung Hom)
\end{itemize}}

\end{frame}

\end{document}
